% TODO make hw stylesheet
\documentclass[10pt]{article}
\usepackage[letterpaper,top=1in,left=1in,right=1in]{geometry}
\usepackage[dvipsnames,svgnames]{xcolor}
\usepackage[shortlabels]{enumitem}
\usepackage[framemethod=TikZ]{mdframed}
\usepackage{amsmath,amssymb,amsthm}
\usepackage[colorlinks]{hyperref}
\usepackage{multicol}
\usepackage{tabularx}
\usepackage{bbm}

\hypersetup{
    colorlinks=true,
    linkcolor=blue,
    filecolor=magenta,
    urlcolor=magenta,
    pdftitle={MAT 344 HW}
}

\usepackage{mathtools}
\usepackage{thmtools}
\usepackage[textsize=footnotesize]{todonotes}
\usepackage{titling}
\usepackage{cancel}

\reversemarginpar
\mdfdefinestyle{mdbluebox}{roundcorner=10pt,innerbottommargin=9pt,
linecolor=blue,backgroundcolor=TealBlue!5,}
\declaretheoremstyle[headfont=\sffamily\bfseries\color{MidnightBlue},
mdframed={style=mdbluebox},]{thmbluebox}

\mdfdefinestyle{mdbloobox}{frametitlefont=\bfseries,innerbottommargin=8pt,
nobreak=true,backgroundcolor=TealBlue!5,linecolor=blue,}
\declaretheoremstyle[headfont=\bfseries\color{MidnightBlue},
mdframed={style=mdbloobox},headpunct={\\[3pt]},postheadspace=0pt,]{thmbloobox}

\mdfdefinestyle{mdredbox}{frametitlefont=\bfseries,innerbottommargin=8pt,
nobreak=true,backgroundcolor=Salmon!5,linecolor=RawSienna,}
\declaretheoremstyle[headfont=\bfseries\color{RawSienna},
mdframed={style=mdredbox},headpunct={\\[3pt]},postheadspace=0pt,]{thmredbox}

\mdfdefinestyle{mdgreenbox}{linecolor=ForestGreen,backgroundcolor=ForestGreen!5,
linewidth=2pt,rightline=false,leftline=true,topline=false,bottomline=false,}
\declaretheoremstyle[headfont=\bfseries\sffamily\color{ForestGreen!70!black},
mdframed={style=mdgreenbox},headpunct={ --- },]{thmgreenbox}

\mdfdefinestyle{mdorangebox}{linecolor=RedOrange,backgroundcolor=RedOrange!5,
linewidth=2pt,rightline=false,leftline=true,topline=false,bottomline=false,}
\declaretheoremstyle[headfont=\bfseries\sffamily\color{RedOrange!70!black},
mdframed={style=mdorangebox},headpunct={ --- },]{thmorangebox}

\mdfdefinestyle{mdblackbox}{linecolor=black,backgroundcolor=RedViolet!5!gray!5,
linewidth=3pt,rightline=false,leftline=true,topline=false,bottomline=false,}
\declaretheoremstyle[mdframed={style=mdblackbox}]{thmblackbox}

\declaretheoremstyle[spaceabove=6pt,spacebelow=6pt]{boldhead}
\declaretheoremstyle[spaceabove=6pt,spacebelow=6pt,headfont=\normalfont\itshape]{italichead}

\declaretheorem[style=boldhead,name=Problem]{problem}
\declaretheorem[style=boldhead,name=Problem,numbered=no]{problem*}
\declaretheorem[style=boldhead,name=Setup]{setup} % for config geo
\declaretheorem[style=boldhead,name=Example,sibling=problem]{example}
\declaretheorem[style=boldhead,name=Theorem,sibling=problem]{theorem}
\declaretheorem[style=boldhead,name=Corollary,sibling=theorem]{corollary}
\declaretheorem[style=boldhead,name=Proposition,sibling=theorem]{prop}
\declaretheorem[style=boldhead,name=Lemma,numberwithin=problem]{lemma}
\declaretheorem[style=boldhead,name=Theorem,numbered=no]{theorem*}
\declaretheorem[style=boldhead,name=Lemma,numbered=no]{lemma*}
\declaretheorem[style=boldhead,name=Claim,numberwithin=problem]{claim}
\declaretheorem[style=boldhead,name=Claim,numbered=no]{claim*}
\declaretheorem[style=boldhead,name=Remark,numberwithin=problem]{remark}
\declaretheorem[style=boldhead,name=Remark,numbered=no]{remark*}
\declaretheorem[style=boldhead,name=Definition,sibling=theorem]{definition}
\declaretheorem[style=boldhead,name=Definition,numbered=no]{definition*}

\providecommand{\ol}{\overline}
\providecommand{\ceil}[1]{\left\lceil #1 \right\rceil}
\providecommand{\floor}[1]{\left\lfloor #1 \right\rfloor}
\newcommand{\dd}{\mathrm{d}}
\providecommand{\ul}{\underline}
\providecommand{\eps}{\varepsilon}
\providecommand{\half}{\frac{1}{2}}
\providecommand{\CC}{\mathbb C}
\providecommand{\FF}{\mathbb F}
\providecommand{\II}{\mathbb I}
\providecommand{\NN}{\mathbb N}
\providecommand{\QQ}{\mathbb Q}
\providecommand{\RR}{\mathbb R}
\providecommand{\ZZ}{\mathbb Z}
\providecommand{\ind}{\mathbbm 1}
\providecommand{\dg}{^\circ}
\providecommand{\ii}{\item}
\providecommand{\setdiff}{\smallsetminus}
\providecommand{\alert}[1]{{\sffamily\textbf{\textcolor{blue}{#1}}}}
\providecommand{\inv}{^{-1}}
\providecommand{\mb}[1]{\mathbf{#1}}
\newcommand{\what}{\widehat}

\newenvironment{soln}{\begin{proof}[Solution]}{\end{proof}}

\providecommand{\pow}{\operatorname{Pow}}
\providecommand{\vocab}[1]{{\textbf{\textcolor{ForestGreen}{#1}}}}
\providecommand{\lcm}{\operatorname{lcm}}
\providecommand{\abs}[1]{\left\lvert #1\right\rvert}
\providecommand{\smabs}[1]{\lvert #1\rvert}
\providecommand{\innerprod}[1]{\langle\!\langle #1\rangle\!\rangle}
\providecommand{\norm}[1]{\left\| #1\right\|}
\providecommand{\smnorm}[1]{\| #1\|}
\providecommand{\gr}{\operatorname{gr}}
\providecommand{\im}{\operatorname{im}}
\providecommand{\paren}[1]{\left(#1\right)}

\usepackage{mathrsfs}

% place title at top right
\setlength{\droptitle}{-8ex}
\pretitle{\begin{flushright}\Large\bfseries}
\posttitle{\par\end{flushright}}
\preauthor{\begin{flushright}}
\postauthor{\end{flushright}}
\predate{\begin{flushright}}
\postdate{\end{flushright}}

\title{MAT 344 HW10}
\author{\large Aditya Pahuja}
\date{\large Due 05/05}
\begin{document}
\maketitle

\begin{problem}
    [p.227\#2]
    Show that the functions $z^n$ for nonnegative integer $n$
    form a normal family in the regions $\abs z < 1$
    and $\abs z > 1$, but not in any region containing
    a point on the unit circle.
\end{problem}

\begin{soln}
    For every point $z$ in the unit disk,
    $\abs{z^n} \le 1$, so by Montel's theorem
    $\mathcal F = \{z^n : n\in\ZZ_{\ge 0}\}$ is normal
    on the open unit disk $\Delta$.

    To show normality on $\Omega = \{z : \abs z > 1\}$,
    let $\{z^{n_k}\}\subseteq\mathcal F$ be a sequence.
    If $\{n_k\}$ is bounded, then some integer $m$ appears infinitely often
    among the $n_k$, so that there is a subsequence
    $z^m$, $z^m$, \dots which trivially converges uniformly to $z^m$.
    Otherwise, if $\{n_k\}$ is unbounded,
    replace $\{n_k\}$ with a subsequence $\{m_k\}$
    which is strictly increasing.
    When $\abs z \ge 1 + \eps$ for some fixed $\eps > 0$,
    we have that for any given $M > 0$,
    if $k > \log_{1+\eps}(M)$, then
    \begin{equation*}
        \abs{z^{m_k}} \ge (1 + \eps)^k > M.
    \end{equation*}
    That is, $\{z^{m_k}\}$ converges uniformly to $\infty$
    on the compact annulus $1 + \frac 1n \le \abs z \le n$
    for each positive integer $n$.
    These annuli exhaust $\{z : \abs z > 1\}$, so we have found
    a subsequence of $\{z^{n_k}\}$ which converges locally uniformly,
    proving that $\mathcal F$ is normal on this region.

    Finally, suppose that $\Omega$ is a domain containing
    a point $a$ on the unit circle.
    Let $U\subseteq\Omega$ be a neighborhood of $a$.
    The sequence $z$, $z^2$, $z^3$, \dots (and any subsequence thereof)
    converges pointwise to $0$ on $U\cap\Delta$,
    $1$ on $U\cap\partial\Delta$, and $\infty$
    on $U - \ol\Delta$; this pointwise limit is not continuous,
    so the convergence cannot have been uniform,
    meaning $\mathcal F$ is not normal on $\Omega$.
\end{soln}

\begin{problem}
    [p.227\#3]
    If $f(z)$ is entire, show that the family of functions
    $\mathcal F = \{f(kz) : k\in\CC\}$ is normal in the annulus
    $r_1 < \abs z < r_2$ if and only if $f$ is a polynomial.
\end{problem}

\begin{soln}
    Suppose $f$ is a polynomial.
    If $f$ is constant then $\mathcal F$ is obviously normal, so suppose
    otherwise and let $\{f(k_nz)\}_{n\in\NN}$ be a sequence in $\mathcal F$.
    If $\{k_n\}$ is unbounded, then pick a subsequence $\{k_{n_j}\}$
    for which $\abs{k_{n_j}}\to\infty$;
    then $\lim_{j\to\infty} f(k_{n_j}z) = \infty$
    for each $z$ in the annulus. This convergence is
    uniform on compact subsets of the annulus because
    for a polynomial $p(z) = a_nz^n + \cdots + a_1z + a_0$,
    $\frac{p(z)}{z^n}$ is arbitrarily close to $a_n$
    as $z$ grows arbitrarily large, so that
    for any $\eps > 0$ there is an $R > 0$ such that if $\abs z > R$ then
    $\abs{p(z)}$ is bounded below by $(\abs{a_n} - \eps)\abs{z^n}$.
    If on the other hand $k_n$ is a bounded sequence, then
    \begin{equation*}
	\sup_{\substack{n\in\NN \\ r_1 < \abs z < r_2}} \abs{f(k_nz)} < \infty,
    \end{equation*}
    which implies the existence of a locally uniformly convergent
    subsequence on the annulus by Montel's theorem.

    Conversely, suppose $\mathcal F$ is normal on the annulus and
    consider the sequence $\{f(nz)\}_{n\in\NN}$.
    Let $s_1,s_2\in\RR$ such that $r_1 < s_1 < s_2 < r_2$;
    then a subsequence of $f(nz)$ converges uniformly on
    $s_1\le\abs z\le s_2$. If $\{f(nz)\}_{n\in\NN}$
    a uniformly bounded sequence, then $f$ is a bounded entire function,
    so it is constant. If this sequence is unbounded,
    by applying normality to some subsequence for which
    $\sup_{r_1 < \abs z < r_2} \abs{f(nz)}$ has limit $\infty$
    shows that this subsequence has a further subsequence
    which converges locally uniformly to $\infty$.
    This implies that $f$ has a pole at $\infty$,
    which for entire functions is equivalent to being a polynomial.
\end{soln}

\pagebreak

\begin{problem}
    [p.227\#4]
    If a family $\mathcal F$ of holomorphic (or meromorphic)
    functions is not normal in $\Omega$,
    show that there exists a point $z_0$ such that
    $\mathcal F$ is not normal in any neighborhood of $z_0$.
\end{problem}

\begin{soln}
    Suppose that for each $z\in\Omega$ there is
    a neighborhood $U_z\subseteq \Omega$ of $z$
    in which $\mathcal F$ is normal.
    Let $K\subseteq\Omega$ be compact
    and $\{f_n\}\subseteq\mathcal F$ a sequence.
    There are finitely many points $z_1$, \dots, $z_m$
    such that the $U_{z_m}$ cover $K$.
    Let $\{f_{n_{1,k}}\}_{k\in\NN}$ be a subsequence of $\{f_n\}$
    which converges locally uniformly on $U_{z_1}$,
    and, for $1 < i \le m$, define $\{f_{n_{i,k}}\}_{k\in\NN}$
    to be a subsequence of $\{f_{n_{i-1,k}}\}_{k\in\NN}$
    which converges locally uniformly on $U_{z_i}$,
    which we can do by the normality assumptions.
    That is, for each $K$ we can find a subsequence
    $\{f^K_n\}_{n\in\NN}$ of $\{f_n\}$ which converges uniformly on $K$.

    Now let $K_1\subseteq K_2\subseteq\cdots\subseteq\Omega$
    be an exhaustion of $\Omega$ by compact sets $K_j$
    and choose $\{f^{K_j}_n\}_{n\in\NN}$
    to be a subsequence of $\{f^{K_{j-1}}_n\}_{n\in\NN}$ for each $j > 1$,
    and consider the subsequence $\{f^{K_n}_n\}_{n\in\NN}$.
    This sequence necessarily converges uniformly on each $K_j$
    because it is a subsequence of $\{f^{K_j}_n\}_{n\in\NN}$
    for each $j$, so we have shown that $\mathcal F$ being
    locally normal implies that it is normal;
    the contrapositive is what we wanted to prove.
\end{soln}










\end{document}
